\scanpage{241}
\mspage{191}

3. \uline{Sur la philosophie des fibres homotopiques.}

\deleted{\unsure{3.}} Le critère d'asphéricité (pour $W_\infty$) que je viens d'obtenir, grâce aux foncteurs $\underline{\mathrm{Ch}}(X)$ (ou au choix, $\underline{\mathrm{Ch}}_\infty(X)$) et ses dérivés $\Omega(X,x)$ et les itérés, enlève tout doute (s'il pouvait en rester un) que les \deleted{\unsure{esp.}} ``espaces de chemins'' que j'ai construit (assez péniblement !) sont bel et bien ``les bons''.

Je voudrais maintenant me rendre compte de façon la plus précise possible \deleted{\unsure{des}} de leur ``canonicité'', dans quel mesure ces constructions sont indépendantes d'un choix (p.~ex. celui entre \unsure{la|le} \unsure{théorie|choix} $\underline{\mathrm{Ch}}$, et $\underline{\mathrm{Ch}}_\infty$).

Tout revient, finalement, à des questions de factorisations d'une flèche quelconque

\[
f:X\longrightarrow Y
\]

de $\mathrm{Cat}$ en

\diagram{\begin{tikzcd}[ampersand replacement=\&]
X\arrow[r,"{\deleted{\unsure{f}}i_{\deleted{\unsure{f}} \unsure{f}}}"]\&\overline X\arrow[r,"{\overline f}"]\&Y,
\end{tikzcd}}

avec

\[
\unsure{i_f}\in W,\quad \overline f\in W\text{-}\mathrm{fib}
\]

(le fait que $\overline f$ puisse $=$ être choisie $W_\infty$-parfaite ne me paraissant pas essentiel ici).

\unsure{C'est} bien de cela qu'il s'agit pour construire
% archive sheet number: 241
